% page 595 of the facsimile (image p-594 of the scan)
% block 165.1 x 246.3 mm ; left margin 36.9 mm
\pgc{15.240mm}{0.000mm}{
\l{\sou{tortikolo.}\cel{1}(\sou{patol.})\cel{2}Doloro che la muskuli di la kolo, qua}
\l{\cel{3}koaktas ad igar olca deviacar de vertikalo. - DEFIS.}
\l{}
\l{\sou{tortugo}. (\sou{zool.)}\cel{2}Genero de repteri, amfibia, quar-peda, qua}
\l{\cel{3}marchas lente, kun kuraso de materio korno-harda (skalio).}
\l{\cel{3}- L.\cel{1}\sou{chelonia}. - EFI.}
\l{}
\l{\sou{tostar}. (\sou{netrans.,}\cel{1}ad)\cel{2}Propozar (a la personi prezenta) drinkor}
\l{\cel{3}\textquotedbl{}por la duro di la saneso\textquotedbl{} (di ulu), akompanante ta propozo}
\l{\cel{3}per diskuraeto. - DEFR.}
\l{}
\l{\sou{tota}. I. De qua nulo deprenesis per separo. -II. (a) Qua nule}
\l{\cel{3}diminutesis. (b) (\sou{metaf.}) Qua ne implikas restrikteso. -}
\l{\cel{3}EFI.}
\l{}
\l{\sou{totalizatoro}. (\sou{tekn.}) Aparato qua donas la sumo di serio de ope-}
\l{\cel{3}raci. -}
\l{}
\l{\sou{toxiko}. Substanco venenoza. - DEFIRS.}
\l{}
\l{\sou{toxikologio}. La cienco qua traktas la toxiki. - DEFIS.}
\l{}
\l{\sou{toxino}. (\sou{kemio biol.}) Veneno komplexa e solvebla, quan produktas}
\l{\cel{3}la funciono di la mikrobi, sive en la organismi vivanta,}
\l{\cel{3}sive en la medii artificala di kultivo. - DEFIS.}
\l{}
\l{\sou{tra.}\cel{3}Irante de ca a ta extremajo, sinse di profundeso.FS.}
\l{}
\l{\sou{trabo.}\cel{2}(\sou{teknol.)}\cel{3}Stangego de ligno, quadratigata, destinata}
\l{\cel{3}a suportor ula parto di edifico (tekto, plankaro, e c.)- ISL.}
\l{}
\l{\sou{trabekulo.}\cel{2}(\sou{histologio)}.\cel{2}Dina prolonguro qua separas su de}
\l{\cel{3}parieto di celulo e salias en kavajo, maxim-multa-kaze per}
\l{\cel{3}kontraktar anastomosi kun la trabekuli vicina. -}
\l{}
\l{\sou{traco}. I. Serio de marki lasita da korpo presita aden altra, per}
\l{\cel{2}qui onu deskovras la direciono ye qua iris ulu, ulo. -}
\l{\cel{3}II. Singla ek ta marki. (\sou{anke metafore}).- III. Marko per qua}
\l{\cel{3}onu deskovras la efiko da ulo, da ulu, en epoko pasinta. -}
\l{\cel{3}- EFI.}
\l{}
\l{\sou{traciono}. (\sou{.teknol.}) (\sou{pri relvoyo}). Ensemblo de la mashini qui}
\l{\cel{3}tiras, tranas, la treni qui transportas la pasajeri o la vari.}
\l{\cel{3}DEFIS.}
\l{}
\l{\sou{tradeskancio.}\cel{1}(\sou{bot.)}\cel{2}Genero de \textquotedbl{}komelinacei\textquotedbl{}, kun flori violea,}
\l{\cel{3}purpurea o blanka, qua vivas en la regioni atmosfero-varma di}
\l{\cel{3}la suda parto di Amerika. - L.\cel{1}\sou{tradescantia}.-}
\l{}
\l{\sou{tradicionar}. (\sou{trans.)}\cel{3}Transmisar vocale la fakti historiala,}
\l{\cel{3}la doktrini religiala, la kustumi. - DEFIRS.}
\l{}
\l{tradukar. (\sou{trans., de, ad)}\cel{3}Pasigar de linguo ad altra linguo.}
\l{\cel{3}- FIS.}
\l{}
\l{trafiko. Importo, frequeso-grado di la paso di la treni, di la}
\l{\cel{3}navi. - EFIS.}
}
